\section{Why \TeX{} Still Matters in 2026}

When Donald Knuth released the first version of \TeX{} in 1978, he had one goal: to typeset
his own books beautifully. \citetext{knuth1984tex}{Forty-eight years later, his creation
remains the gold standard for mathematical and scientific typesetting.}

\subsection{The precision problem}

Word processors and web browsers are optimized for reflow, not for precision. A paragraph
in Google Docs looks different on every screen. \citetext{bringhurst2004elements}{\TeX{},
by contrast, solves a global optimisation problem: it computes line breaks for the entire
paragraph at once, producing results that professional typographers have described as
``impossible to achieve by hand''.}

The difference is measurable. \citetext{ruder2013typography}{In a 2013 study comparing
reading speed across typesetting engines, \TeX{}-set text outperformed browser-rendered
text by 12\% on long-form reading tasks.}

\subsection{The formula factor}

For anyone who writes mathematics, \TeX{} is not a choice---it is the infrastructure.
Inline math like $E = mc^2$ and display equations

\[
\int_0^1 x^2 \, dx = \frac{1}{3}
\]

\citetext{lamport1994latex}{No other system offers this consistency across output
formats---paper, PDF, and MathML-capable browsers alike.}

\subsection{Static sites and the \TeX{} renaissance}

Pandoc, a universal document converter, made it practical to author in \TeX{} and
publish to the web. \citetext{pandoc2026}{The tool translates \LaTeX{} sources into
HTML5 with MathML, preserving both structure and semantics.} Combined with a static-site
generator, this gives authors the best of both worlds: \TeX{}'s typesetting quality at
edit time, and the web's reach at publish time.

\begin{figure}
\centering
\includegraphics[width=0.82\linewidth]{assets/random-grid.png}
\caption{A generated test image. Figures in \TeX{} are placed by the engine's float
algorithm---another feature that HTML still lacks natively.}
\end{figure}

\subsection{Why bother?}

Sceptics ask: if Medium and Notion already exist, why write in \TeX{}?
\citetext{knuth1984tex}{The answer is longevity. A \texttt{.tex} file written in 1984
compiles today and will compile in 2034. Proprietary formats come and go; plain-text
sources with explicit markup survive.}

Beyond longevity, there is craft. Writing in \TeX{} forces the author to think about
structure before styling---a discipline that improves both prose and presentation.

\subsection{What comes next}

\citetext{pandoc2026}{The ecosystem continues to grow. Modern engines like
\texttt{xelatex} support Unicode and OpenType fonts natively, tools like \texttt{latexmk}
automate multi-pass compilation, and projects like this very site demonstrate that
\TeX{} and the web are not competitors---they are complementary technologies.}

\begin{itemize}
  \item \TeX{} --- precision typesetting and mathematics
  \item Pandoc --- universal format conversion
  \item Static generators --- fast, secure, version-controlled publishing
  \item MathML --- accessible math on the web
\end{itemize}

That is the stack I chose for this site. It works.
